\section{Introduction}\label{sec:introduction}

Haiden, Katzarkov, Kontsevich and Pandit associate a mass and a metric
phase flow to a complex of quiver representations
\cite[\S\S4.2 and 5.4]{hkkp26}. Their construction uses a coefficient
of a superconnection exponential. Acyclic summands should disappear when
the differential becomes large in the evolving metric, leaving the
cohomology representation. Question~5.18 of \cite{hkkp26} asks whether
the mass is nonincreasing along the flow. We give a negative answer for
the phase determined by their positivity condition and degree branches.
We then prove that quantitative acyclic cancellation can hold despite
this failure.

The distinction between dissipation and cancellation is essential.
The first concerns the sign of a derivative at a finite metric. The
second requires the differential to escape to infinity under the
actual equation, together with estimates on the coupled arrows.
Prescribing a large differential and computing a static limit does not
establish that a trajectory enters this regime.

\subsection{The complexes and their phase flow}
Fix a quiver with two vertices and $k$ arrows from the first vertex to
the second. We consider the two-term complex
\begin{equation}\label{eq:complex}
 (V\xrightarrow{(A_1,\ldots,A_k)}W)
 \xrightarrow{(D,0)}(U\longrightarrow0),
\end{equation}
in degrees zero and one. The chain condition is automatic. All vector
spaces are finite dimensional over $\CC$; the Hermitian metrics vary,
whereas the complex maps are fixed. The central charges $z_0,z_1$ have
strictly positive imaginary parts.

Section~\ref{sec:coefficients} derives the coefficient
$\Omega=\Omega_V\oplus\Omega_U\oplus\Omega_W$ explicitly.
On a block of degree $q$, the matrix $T=(-1)^q\Omega$ has strictly
positive imaginary part. If $T=U_T|T|$ is its polar decomposition, set
\begin{equation}\label{eq:phaseintro}
 \Phi=\Arg U_T-\pi q\id,\qquad
 m(h)=\norm{\Omega_V}_1+\norm{\Omega_U}_1+\norm{\Omega_W}_1,
 \qquad h^{-1}\dot h=-\Phi .
\end{equation}
Here $\Arg U_T$ has spectrum in $(0,\pi)$ and
$\norm{M}_1=\tr(M^\dagger M)^{1/2}$ is the trace norm.
For a single matrix, an unadorned norm denotes the operator norm.
Squared lengths of arrow tuples are sums of squared Hilbert--Schmidt
norms, denoted by the subscript $\HS$ on each arrow.
This is the unique phase in the specified branch for which
$\Omega e^{-\ii\Phi}$ is positive definite. In particular, it realizes
the positivity definition of \cite[\S4.2]{hkkp26}.

The coefficient need not be normal, contrary to the normality assertion
in \cite[Lemma~5.16]{hkkp26}. We exhibit a small example in
Section~\ref{sec:coefficients}. This issue is separate from the mass
counterexample: positivity still defines a smooth Hermitian phase, and
the metric equation exists for every finite positive time.
Throughout the paper, the phase means \eqref{eq:phaseintro}, not an
analytic logarithm of a nonnormal coefficient.

\subsection{Mass increase and cancellation}
The following two statements describe different properties of the same
equation.

\begin{theorem}\label{thm:introincrease}
Let $V=U=\CC^2$, $W=\CC$, $k=1$, and take standard initial metrics with
\[
 D=\begin{pmatrix}1/5&0\\0&2\end{pmatrix},\qquad
 A=\begin{pmatrix}1&2\end{pmatrix},\qquad
 z_0=1+\ii\epsilon,\quad z_1=4+\ii\epsilon .
\]
For every sufficiently small $\epsilon>0$, the initial derivative of
$m$ along \eqref{eq:phaseintro} is strictly positive.
\end{theorem}

The proof in Section~\ref{sec:increase} reduces the sign to a finite
real calculation with a positive interval bound. The endpoint
$\epsilon=0$ is used only as a one-sided limit; it is not an admissible
central charge. No particular positive upper endpoint for $\epsilon$
is needed for the theorem.

\begin{theorem}\label{thm:introcancellation}
Suppose that $D:V\to U$ is invertible, $\dim V=\dim U=n\geq1$ and
$\dim W=m_0\geq1$. Allow arbitrary $k$ arrows, arbitrary positive initial
metrics, and arbitrary $z_0,z_1$ in the strict upper half-plane.
Put $a=\arg z_1$. The phase flow exists for all $t\geq0$. In moving
orthonormal frames, with $X_j=A_jD^{-1}$, there are constants $C,C'>0$
depending on the initial data and charges such that
\begin{align}
 \sum_j\norm{X_j(t)}_{\HS}^2
 &\leq e^{-at}\sum_j\norm{X_j(0)}_{\HS}^2,\label{eq:introX}\\
 \lambda_{\min}(D^\dagger D)(t)&\geq at-C,\label{eq:introlambda}\\
 0\leq m(h(t))-m_0|z_1|&\leq C'e^{-at}.\label{eq:intromass}
\end{align}
The acyclic coefficient blocks tend to zero and
$\Omega_W\to z_1\id_W$ at the same exponential rate in the evolving
trace norms. In fixed coordinates on $W$, the metric
$e^{at}h_W(t)$ converges exponentially to a positive definite limit.
\end{theorem}

The proof uses the exact contraction equation for $A_jD^{-1}$.
A weighted polar estimate then forces the least singular value of $D$
to grow. Neither step assumes that the mass decreases.
No simultaneous diagonalization of $D^\dagger D$ and
$\sum_j A_j^\dagger A_j$ is imposed.

\subsection{Cancellation with arbitrary surviving cohomology}
When $\ker D$ is nonzero, the cohomology has both a source and a target.
The ratio equation acquires an additional term, which can still be
controlled without a stability or balance assumption on the surviving
representation. Let $H=\CC^{n_s}$, $W=\CC^{n_w}$, $U=\CC^r$ and
$V=H\oplus\CC^r$, with $n_s,n_w,r\geq1$. In initial adapted
orthonormal frames take arbitrary arrows $B_j:H\to W$ and set
\begin{equation}\label{eq:mixedinitial}
 D(0)=[0,F_0],\qquad A_j(0)=[B_j,T_j],\qquad
 F_0\in\GL(r,\CC),\quad \sum_j\norm{T_j}_{\HS}^2\leq M^2.
\end{equation}
Fix any charges $z_0,z_1$ in the strict upper half-plane and put
$a=\arg z_1$. Only the least singular value of $F_0$ will be constrained.
Its other singular values may be arbitrarily large.

Along the flow use the orthogonal splitting
$V=\ker D\oplus(\ker D)^\perp$ and write $D=[0,F]$, $A_j=[H_j,T_j]$.
Let $\Omega_H$ be the ordinary degree-zero coefficient of the
instantaneous surviving arrows $H_j$, and let $m_H$ be its trace-norm
mass. Its inclusion into the full graded space places zero in the
acyclic blocks. Put
\[
 q=\sum_j\norm{H_j}_{\HS}^2,\qquad
 s=\sum_j\norm{T_jF^{-1}}_{\HS}^2,\qquad
 \lambda=\sigma_{\min}(F)^2,\quad\lambda_0=\lambda(0),
\]
and set
\[
 Q=\tfrac12\bigl(m_H(0)-n_s\Rea z_0+n_w\Rea z_1\bigr)+1.
\]
The elementary mass bound in Section~\ref{sec:mixed} gives
$Q\geq q(0)+1$.

\begin{theorem}\label{thm:intromixed}
There are $\Lambda_*(Q,M,z_0,z_1)$ and
$C(n_s,n_w,r,Q,M,z_0,z_1)$ such that every initial complex
\eqref{eq:mixedinitial} with
$\sigma_{\min}(F_0)\geq\Lambda\geq\Lambda_*$ has a global phase flow
satisfying $q(t)\leq Q-3/4$ and
\begin{equation}\label{eq:intromixedratios}
 s(t)\leq M^2\Lambda^{-2}e^{-at/2},\qquad
 \lambda(t)\geq\max\{\lambda_0,at/4-C\}.
\end{equation}
Moreover,
\begin{equation}\label{eq:intromixed}
 \norm{\Omega(t)-(\Omega_H(t)\oplus0_{\ac})}_1+
 |m(h(t))-m_H(t)|\leq C\Lambda^{-2}e^{-at/8}.
\end{equation}
The constants are independent of arrow count, the surviving stabilizer
and the condition number of $F_0$. They can be chosen uniformly when
the surviving dimensions, acyclic rank, charges and total initial
arrow-norm bounds are fixed. No balance or stability hypothesis is required.
\end{theorem}

Section~\ref{sec:mixed} gives explicit constants and a threshold
independent of acyclic rank. The charges are fixed in these uniformity
statements; no uniformity as their imaginary parts vanish is asserted.
For example, acyclic singular values $\Lambda$ and $R\Lambda$ are
permitted with the same constants for every $R\geq1$.

The theorem compares with evolving cohomology. It does not assert that
cohomology itself approaches balance. The proof uses the ordinary
cohomology mass to bound $q$, and a polar estimate using only the
nonsingular cohomology block to control the additional coupling.
A separate weighted phase estimate forces growth of the acyclic
differential. These arguments make no use of full-complex mass
dissipation or of a positive orbit Hessian gap.

\subsection{Relaxation from finite cohomology energy}
Now suppose $n_s=n_w=n$ and $z_0=-1+\ii$, $z_1=1+\ii$.
Define the moment-map defects
\[
 \mu_s=\sum_jH_j^\dagger H_j-\id,\qquad
 \mu_w=\sum_jH_jH_j^\dagger-\id,\qquad
 \norm\mu^2=\norm{\mu_s}_{\HS}^2+\norm{\mu_w}_{\HS}^2.
\]
In fixed holomorphic coordinates let $h_s,h_w$ be the induced metrics
on $\ker D,W$, both multiplied by $e^{\pi t/2}$.
Using the initial orthonormal determinant volume forms, define the
cohomology Kempf--Ness function on all positive metric pairs by
\[
 \mathcal K=q+\log\det h_s-\log\det h_w.
\]

\begin{theorem}\label{thm:relaxation}
Under the hypotheses of Theorem~\ref{thm:intromixed}, with these equal
dimensions and charges, suppose
$\mathcal K_*:=\inf\mathcal K> -\infty$ and put
$E_0=\mathcal K(0)-\mathcal K_*$. Set $J_Q=1+(Q-1)^2$.
Then
\begin{equation}\label{eq:intromixedenergy}
 \begin{gathered}
 0\leq\mathcal K(t)-\mathcal K_*,\\
 \mathcal K(t)-\mathcal K_*+
 \frac1{2J_Q}\int_0^t\norm{\mu(u)}^2\,du
 \leq E_0+C_{\mathrm{en}}(Q,M)\Lambda^{-4}.
 \end{gathered}
\end{equation}
Both moment-map defects tend to zero and $m(h(t))\to2n$.
The infimum need not be attained, and initial balance is not required.
\end{theorem}

A balanced metric pair attains a finite minimum by geodesic convexity,
but its value depends on the chosen determinant normalization.
Section~\ref{sec:mixed} also gives a finite, unattained example.
A uniform energy bound across initial metrics retains a bound on $E_0$.
Neither this theorem nor the cancellation estimate asserts a uniform
final relaxation rate or convergence of normalized surviving metrics.

The initially balanced case has a particularly simple normalization.
Suppose
\begin{equation}\label{eq:initialbalance}
 \sum_jB_j^\dagger B_j=\sum_jB_jB_j^\dagger=\id_n.
\end{equation}
Then $q(0)=n$, $m_H(0)=2n$, $Q=2n+1$ and
$\mathcal K_* =\mathcal K(0)=n$.

\begin{corollary}\label{cor:balanced}
For \eqref{eq:initialbalance}, the preceding conclusions hold with
threshold $\Lambda_*(n,M)$ and constants $C(n,r,M)$, uniformly over
all balanced tuples. In particular $q(t)\leq2n+1/4$, the cancellation
rate in \eqref{eq:intromixed} is $\pi/32$, and
\eqref{eq:intromixedenergy} holds with $E_0=0$ and $J_Q=1+4n^2$.
\end{corollary}

No individual arrow is required to be invertible, and the surviving
representation need not be irreducible.

The following family makes the stabilizer uniformity explicit. Let
$n=r=2$, $k=3$, and put
\begin{equation}\label{eq:balancedfamily}
 B_1=c\id,\qquad B_2=c\sigma_z,\qquad B_3=c\epsilon\sigma_x,\qquad
 c=(2+\epsilon^2)^{-1/2},\quad 0\leq\epsilon\leq1,
\end{equation}
where
$\sigma_z=\diag(1,-1)$ and
$\sigma_x=\left(\begin{smallmatrix}0&1\\1&0\end{smallmatrix}\right)$.

\begin{corollary}\label{cor:pauli}
For \eqref{eq:balancedfamily}, $F_0=\Lambda\id_2$ and every mixing
tuple with $\sum_j\norm{T_j}_{\HS}^2\leq M^2$, the conclusions of
Corollary~\ref{cor:balanced} hold for $\Lambda\geq\Lambda_*(2,M)$,
uniformly in $\epsilon\in[0,1]$. The cohomology automorphism algebra
is scalar for $\epsilon>0$ and diagonal for $\epsilon=0$.
\end{corollary}

\begin{proof}
Both Gram sums in \eqref{eq:initialbalance} are $\id_2$.
Since $B_1=c\id$ is invertible, an automorphism is a common source
and target matrix commuting with $\sigma_z$ and $\epsilon\sigma_x$.
The first condition makes it diagonal, and the second makes it scalar
when $\epsilon>0$. The estimates follow from Corollary~\ref{cor:balanced}.
\end{proof}

The parameter $\epsilon$ is fixed along each trajectory, but may
approach zero independently of $\Lambda$ across the family.

\subsection{Static regularity and potentials}
The coefficient formulas yield an exact complex primitive in metric
directions. For an arbitrary fixed differential with least acyclic singular
value $\Lambda$, the relative primitive and its bounded-arrow-coordinate
derivatives have error $O(\Lambda^{-2})$. The mass has a $C^1$ limit with
first-derivative error $O(\Lambda^{-1})$. These bounds do not give arbitrary
derivative estimates for the polar positive factor or its phase.

On the balanced family \eqref{eq:balancedfamily}, with fixed nonzero
mixing, the mass converges in $C^2$ but its Hessian has no uniform
positive H\"older bound. A shrinking-scale profile also implies divergence
of third derivatives at points approaching $\epsilon=0$, and of the
fourth derivative at $\epsilon=0$. Thus the static cancellation limit
does not have unrestricted smooth parameter dependence.
The real phase combination of the primitive also has a negative
ambient Levi direction. Section~\ref{sec:geometry} therefore separates
the positive limiting cohomology potential from a more restricted
finite-scale positive construction. In that construction the quotient
comparison is an exact scalar isometry. It is not a metric on the
full space of arbitrary mixed complexes.

\subsection{Relation to earlier work}
In the zero-arrow case, the differential rescaling and cancellation
mechanism already appear in \cite[Example~4.1]{hkkp26}.
The motivating large-differential discussion in \cite[\S5.4]{hkkp26}
is a static comparison. Our results concern the actual coupled phase
equation and prove that its solutions enter the required regime.
The ordinary quiver metric flow studied in \cite[\S5]{hkkp20}
has a different vector field; its asymptotic results do not by themselves
apply to \eqref{eq:phaseintro}.

Superconnection localization to cohomology and heat-transgression
primitives have classical precedents. Bismut--Lott
\cite[Theorems~2.13 and 2.16]{bl95} treat cohomology localization and
transgression under a prescribed graded metric rescaling.
Paradan--Vergne \cite[\S3.2, Proposition~4 in the preprint]{pv09} give heat-decay
estimates away from the noninvertibility locus for a prescribed
superconnection rescaling. We use elementary finite-dimensional
instances of these analytic ideas. The additional estimates here
control the noncommuting polar vector field, including coupling to
a surviving representation without a positive orbit gap.
The quotient constructions in the last section use classical
quiver and K\"ahler reduction \cite{king94,sj95,sl91}; their role is
to identify precisely the geometry supplied by the explicit primitive.

Sections~\ref{sec:coefficients}--\ref{sec:polar} establish the coefficient
and polar estimates. Section~\ref{sec:invertible} proves invertible-source
cancellation, Section~\ref{sec:increase} proves mass increase, and
Section~\ref{sec:mixed} treats surviving source cohomology.
Sections~\ref{sec:static}--\ref{sec:geometry} concern static regularity
and potentials.
